\subsection{Coleta completa e consolidação do dataset}

A coleta da Sprint 2 utiliza a janela de 1 de outubro de 2025 a 30 de setembro
de 2026, semente 42 e meta mínima de 300 repositórios após os filtros de inclusão.
O comando de coleta permite reaproveitar um snapshot versionado de candidatos;
nesse caso, sua ordem e consultas de origem são preservadas. Uma nova busca pode
ampliar as faixas de estrelas, com subdivisão para respeitar o teto da Search API.
Não se considera a amostra final congelada antes da coleta e validação completas.

O pipeline utiliza Python e \texttt{requests}, sem bibliotecas que encapsulem a
API do GitHub. São consultados os endpoints de workflows, releases, comparações
entre tags, commits, contribuidores e execuções do GitHub Actions. As tags podem
ser recuperadas em lotes por GraphQL. Releases e comparações são paginadas; a
data de autoria do commit é utilizada nos cálculos de lead time e nas datas das
tags. Os runs são restritos ao default branch, ao evento push e à janela do estudo.
Consultas mensais com pelo menos mil resultados são subdivididas; intervalos
diários ainda saturados são registrados como incompletos e impedem o congelamento.

O cache em disco permite reutilizar respostas, enquanto checkpoints registram o
avanço da seleção e as coletas completas por repositório. O cliente respeita os
limites indicados nos cabeçalhos e aplica novas tentativas com espera crescente
em falhas temporárias. Erros transitórios não constituem evidência de exclusão.
Comparações com erro 404 são contabilizadas e excluídas do lead time; datas
inválidas de commits também são registradas e excluídas das medianas.

A execução completa terminou em 15 de outubro de 2026, às 05:15 UTC. O ambiente
de execução utilizou Python 3.12.14. A instrumentação parcial, iniciada após o
começo da coleta, registrou 25.634 respostas HTTP (25.624 com status 200 e dez
com status 404) e três esperas, somando três segundos solicitados. Esses valores
não incluem acertos de cache nem representam todas as requisições desde a seleção
inicial. As esperas registradas não comprovam, isoladamente, ocorrência de rate limit.
As 311 fontes congeladas são distribuídas em arquivo ZIP, com verificação SHA-256;
a restauração e o recálculo sem acesso à API reproduziram exatamente o dataset.

O dataset reutiliza as funções de cálculo do grupo, mantém uma linha por
repositório e registra métricas, versão da heurística e contagens de ignorados
e censurados. O CFR de CI considera falhas sobre falhas mais sucessos. O CFR por
release observa uma corretiva posterior em até sete dias, excluindo do denominador
as releases das sete últimas datas da janela. Episódios de recuperação não
encerrados são mantidos como censurados. Valores ausentes não são substituídos
por zero. Tipos complementares só são incorporados com revisão declarada e
proveniência separada do consenso.
O consenso somente é incorporado quando seu manifesto confirma a origem da
rotulagem e o vínculo, por hashes, com a amostra congelada. O consenso legado
do piloto sem esse vínculo não é considerado validação da nova coleta. A avaliação
definitiva da amostra-ouro utiliza as três planilhas com origem humana declarada
pelos integrantes, rigorosamente vinculadas à população congelada. O consenso
fornece o tipo e a validação para os 60 projetos sorteados, e a classificação
complementar dos outros 240 projetos foi revisada pelos integrantes após sugestões
automáticas e integrada com proveniência e links de evidência, totalizando os 300
projetos com tipo. O importador verifica cobertura exata, categorias, campos
obrigatórios, formato das URLs e datas, hashes e metadados do manifesto. Ele não
comprova automaticamente a realização da revisão humana ou a disponibilidade das
páginas externas. Arquivos já atestados não são sobrescritos durante a conferência.

\IfFileExists{tabelas/coleta_s02.tex}{A execução concluída reuniu 300 repositórios, na janela de 2025-10-01 a 2026-09-30. Foram contabilizados 11453 releases principais, 291405 runs válidos e 22871 runs ignorados. Houve 9 comparações com erro 404 e 561 relações release/commit com data inválida. No CFR (b), 274 releases foram censuradas. Dos 7861 episódios de falha, 90 permaneceram censurados. Foram registrados 56 intervalos saturados (incluindo subdivisões), com 0 incompletos. O tipo do projeto permanece ausente em 300 repositórios não cobertos pelo consenso, sem imputação automática. O CFR (b) ficou indisponível em 7 repositórios por evidência textual insuficiente.
}{
A coleta completa e seus números ainda aguardam o fechamento dos artefatos da S02.
}
